\addcontentsline{toc}{chapter}{Motivation}
\chapter*{Motivation}

Robotics is a field that is experiencing a huge impulse currently. And more specifically, Open Source robotics is growing rapidly. 
\\

The idea of creating common software, of aiding other investigators to easily replicate the work already done and continue from that point instead of "reinventing the wheel" started in the past century. In 1998 the Open Source Initiative (OSI)\cite{osi} was formed, and its definition is recognized as the standard. According to them, "Open source does not just mean access to the source code. The distribution terms of open-source software must comply with the following criteria: Free Redistribution, Source Code, Derived Works, Integrity of The Author's Source Code, No Discrimination Against Persons or Groups, No Discrimination Against Fields of Endeavor, Distribution of License, License Must Not Be Specific to a Product, License Must Not Restrict Other Software,  License Must Be Technology-Neutral"\cite{osi_def}. 
\\

It is my believe that Open Source is critical in the development of new knowledge and not only in the Software field, but in all the science and technical disciplines. This is how the idea of increasing that common pool of tools with another one that might be useful appeared. 
\\
\begin{wrapfigure}{r}{0.3\textwidth}
	\centering
    \includegraphics[width=0.25\textwidth]{img/intro/open_source.eps}
	\caption[Open Source Initiative Logo]{Open Source Initiative Logo}
\end{wrapfigure}

It is noticeable the mark that Open Source has throughout the project. From the libraries being used in it to the Robotic Operating System of which this project is but another package. All of them are Open Source. Also, the tutorials, examples and different web-pages with useful comments and aids for those who are learning has been critical in this project's development. Nothing of this could have been possible without the idea of Open Source code. 
\\

Apart from the want of making useful Open Source code, there remain another motivation unexplained. That is, the specific subject of this project. 
\\

Previously, I have developed other projects involving designing and constructing wheeled robots. It was then the idea of creating different packages that glued together could suit almost every robot appeared. One of the most interesting fields that can be used in virtually all types of robots is Computer Vision. 
\\

Computer vision has experience an important improvement in the last years through the upgrade of the hardware and software that compose it. The hardware such as acquisition elements (cameras, depth sensors, etc) and computing elements (PCs or other programmable devices) have experienced a rapid advance in the past years. It allowed to process more data that is now obtained more accurately and with less noise. This increase in the computing power of the equipment created a possibility of introducing more complex libraries and frameworks and even operating systems. 
\\

But Computer Vision is an ample area, why choosing object recognition? The project being developed is intended to be useful, not just a mere hobby. It occurred to me that apart from being useful in the interaction with robots and cognitive robots more particularly, this software could be used with visually impaired people, as an example. Mainly, with people that had a sudden loss of vision and that is not used to detect the objects by its shape, or cannot differentiate them by the contour. 

